\documentclass[10pt,twocolumn]{article}

\usepackage[margin=0.75in]{geometry}
\usepackage{times}
\usepackage{amsmath}
\usepackage{graphicx}
\usepackage{booktabs}
\usepackage{caption}
\usepackage{hyperref}
\usepackage{xcolor}

\title{\vspace{-1em}Early Encrypted-Traffic QoS Classification and Real-Time Enforcement:\\
A Machine-Learning-Driven Router Testbed}

% TODO(author): replace with your name(s) and affiliation before submission.
\author{Author Name(s)\\ \small Affiliation \\ \small \texttt{email@example.org}}
\date{}

\begin{document}
\maketitle

\begin{abstract}
Modern access-network traffic is almost entirely encrypted, which prevents
deep packet inspection from informing router-level Quality-of-Service (QoS)
decisions. We present an end-to-end system that classifies flows into four
QoS classes -- \emph{Delay-Sensitive}, \emph{Video-Streaming},
\emph{Bulk-Download}, and \emph{Web-Browsing} -- using only the Sequence of
Packet Lengths and Times (SPLT) of the first 10 packets of a flow, and
drives Linux traffic control (\texttt{tc}) and \texttt{iptables} to enforce
differentiated treatment in real time. The classifier is trained on
429{,}597 real, labeled network flows and achieves a macro F1-score of
0.871 (accuracy 90.8\%) using a depth-bounded Random Forest sized for
router deployment (236\,MB). A live sniffer built on the same feature
pipeline classifies flows with a measured end-to-end processing latency of
28.9\,ms on average (\(n=8\), unoptimized development hardware). In a
Docker-based router-in-the-middle pilot testbed under a deliberately
constrained 5\,Mbit/s bottleneck link, enabling the QoS mechanism reduced
mean queuing delay for Delay-Sensitive traffic from 1{,}532.8\,ms to
0.095\,ms and for Web-Browsing traffic from 2{,}590.8\,ms to 0.144\,ms,
while Video-Streaming and Bulk-Download absorbed the resulting congestion
as intended by the priority scheme. We further report a methodologically
important negative finding: synthetic traffic generated by D-ITG does not
reproduce the packet-timing signature of the real encrypted applications
our classifier was trained on, causing most synthetic flows to default to
the majority class. We discuss this synthetic-traffic domain gap and its
implications for evaluating early-classification systems, and directly
test it by replaying genuine packet captures through the live classifier
with \texttt{tcpreplay}, which produces qualitatively healthier,
non-degenerate predictions than synthetic traffic did. We additionally
cross-validate our \texttt{tc} enforcement mechanism with \texttt{iperf3},
an independent bandwidth-measurement tool sharing no code with our own
pipeline, confirming both that the testbed's unshaped capacity is
effectively unconstrained and that our HTB configuration delivers the
rate it claims to within 5\%.
\end{abstract}

\section{Introduction}

Internet service providers and home routers have traditionally relied on
deep packet inspection (DPI) or port-based heuristics to identify
application traffic and apply differentiated Quality-of-Service (QoS)
treatment: low latency for voice and gaming, guaranteed throughput for
video, and best-effort scheduling for bulk transfers. The near-universal
adoption of TLS 1.3 and QUIC has eroded the visibility DPI depends on,
motivating a large body of work on machine-learning-based traffic
classification from statistical or sequence-based flow features that
survive encryption~\cite{anderson2017deciphering}.

Two practical constraints separate a research classifier from a
deployable QoS system. First, the classification decision must be made
\emph{early} -- a router cannot wait for a flow to finish before deciding
how to schedule it. Second, the decision must be \emph{actioned}: a
correct label is only useful if it changes how packets are queued and
forwarded. Much of the traffic-classification literature addresses the
first constraint (feature engineering, model accuracy) without closing
the loop on the second (real-time enforcement).

This paper presents a complete pipeline that does both. We (1) remap a
large corpus of nDPI-labeled real network flows onto a 4-class QoS
taxonomy motivated by router scheduling policy rather than by application
identity; (2) train and size a tree-based classifier that uses only the
first 10 packets of each flow's SPLT sequence, matching what a router can
observe before it must decide; (3) implement a live sniffer that performs
this classification in real time and drives a Linux \texttt{tc}
hierarchical-token-bucket (HTB) scheduler via dynamically installed
\texttt{iptables} firewall-mark rules; and (4) build and exercise a
Docker-based router-in-the-middle testbed with synthetic traffic from
D-ITG to measure classification accuracy, router processing overhead, and
the resulting change in per-class delay, jitter, throughput, and loss.

Throughout, we report only measurements we actually collected rather than
representative or illustrative figures. Where an experiment did not work
as intended -- specifically, D-ITG's synthetic traffic does not resemble
the real encrypted applications the classifier was trained on -- we report
that as a finding rather than omitting it, and we isolate the QoS
enforcement mechanism's effectiveness from classifier accuracy using a
ground-truth (oracle) flow-labeling control, so that the two claims
(``the classifier works on real traffic'' and ``the enforcement mechanism
works given a label'') each rest on their own evidence.

\subsection*{Contributions}
\begin{itemize}
  \item An empirically grounded, documented mapping from 337 raw nDPI
  application labels (429{,}597 real flows, a university-network capture)
  to a 4-class router-scheduling taxonomy, requiring only 6 explicit
  application-name overrides beyond nDPI's own category field.
  \item A size/accuracy engineering study showing that an unconstrained
  Random Forest reaching 0.886 macro F1 pickles to 1.3\,GB -- impractical
  for router deployment -- while a depth-bounded variant reaches 0.871
  macro F1 at 236\,MB.
  \item A real-time sniffer and \texttt{tc}/\texttt{iptables} enforcement
  layer sharing a single feature-extraction implementation with the
  offline training pipeline, eliminating train/serve skew, with measured
  end-to-end processing latency.
  \item A router-in-the-middle pilot testbed (plain Docker networking, in
  lieu of a full Containernet deployment that requires host root) showing
  real, measured delay/jitter/throughput/loss changes when the QoS
  mechanism is enabled versus a flat rate-limited baseline.
  \item A documented negative result on the domain gap between D-ITG
  synthetic traffic and real encrypted application traffic for
  early-classification evaluation, directly tested (not just asserted) by
  replaying genuine packet captures through the live classifier with
  \texttt{tcpreplay}, and cross-validated against an independent
  bandwidth-measurement tool (\texttt{iperf3}) sharing no code with our
  own enforcement or evaluation pipeline.
\end{itemize}

\section{Related Work}

\textbf{Encrypted traffic classification.} Anderson and
McGrew~\cite{anderson2017deciphering} established that sequences of
packet lengths and inter-arrival times (SPLT), combined with TLS metadata,
retain enough signal to classify encrypted flows without payload
inspection. Follow-on work has explored deep sequence models and
statistical flow features; we adopt the SPLT formulation directly because
it is computable incrementally, packet-by-packet, which is what a router
needs.

\textbf{Deep packet inspection at scale.} nDPI~\cite{ndpi} and the
NFStream metering framework built on it provide the ground-truth
application labels our dataset inherits, obtained via full DPI
(confidence level 6, i.e. a complete protocol match) rather than
port-based heuristics.

\textbf{Tree-based classifiers.} Random Forests~\cite{breiman2001random}
and gradient-boosted trees (LightGBM~\cite{ke2017lightgbm}) are the
standard choice for flow classification given their robustness to
heterogeneous, non-normalized numeric features and their comparatively
low inference cost relative to deep networks, which matters directly for
a latency-bounded router deployment.

\textbf{Linux traffic control.} The HTB queuing discipline together with
\texttt{iptables}/\texttt{netfilter} firewall marks is the standard
open-source mechanism for class-based bandwidth allocation on Linux
routers, and is what commodity SOHO router firmware (e.g. OpenWrt) uses
for QoS. We use it directly rather than a simulated scheduler so that our
enforcement measurements reflect real kernel behavior.

\textbf{Emulated testbeds.} Mininet~\cite{lantz2010network} provides
lightweight network emulation via Linux network namespaces; Containernet~\cite{containernet}
extends it to use Docker containers as hosts, allowing each emulated node
to run arbitrary software stacks (here, our classifier and D-ITG). D-ITG
(Distributed Internet Traffic Generator)~\cite{avallone2004ditg}
generates traffic with configurable packet-size and inter-departure-time
distributions and reports per-flow delay, jitter, and loss on decode,
which we use as our QoS-effect measurement instrument.

\textbf{Real-traffic replay and independent measurement.}
\texttt{tcpreplay}~\cite{tcpreplay}, together with its companion
\texttt{tcprewrite}, replays previously captured packets onto a live
interface at controlled speed, optionally rewriting addresses -- we use
it to inject genuine (non-synthetic) traffic into the testbed, directly
testing whether our classifier's behavior differs on real versus D-ITG
traffic. \texttt{iperf3}~\cite{iperf3} is the standard tool for measuring
achievable TCP/UDP throughput between two endpoints; we use it as a
bandwidth ground truth that shares no code with either D-ITG or our own
\texttt{tc} configuration, to independently sanity-check both the
testbed's unshaped capacity and our HTB rate enforcement.

\section{Methodology}

\subsection{Dataset and QoS Taxonomy}
\label{sec:taxonomy}

The underlying dataset and its data-preparation pipeline are those of the
tutorial this work extends~\cite{pekar2025tutorial}: 1{,}142{,}230 flows
captured on a university network with NFStream
(\texttt{accounting\_mode=2}, \texttt{splt\_analysis=25}, TCP/UDP only),
each carrying nDPI's \texttt{application\_name} and
\texttt{application\_category\_name} labels. Following the same cleaning
pipeline, we keep only flows with a full-DPI match
(\texttt{application\_is\_guessed=0}, \texttt{application\_confidence=6}),
a known application name, and at least 10 packets -- the last criterion
coincides exactly with our 10-packet feature window, so every flow we
train on has genuine (non-padded) values in all 10 SPLT positions. This
leaves 429{,}597 flows spanning 337 distinct application labels and 28
nDPI categories.

We map each flow to one of 4 QoS classes using nDPI's
\texttt{application\_category\_name} as the primary signal (e.g. Media,
Video, Streaming, Music~\(\rightarrow\) Video-Streaming; Collaborative,
VoIP, RemoteAccess, Network~\(\rightarrow\) Delay-Sensitive; Download,
SoftwareUpdate, Cloud~\(\rightarrow\) Bulk-Download; everything else
defaults to Web-Browsing), with 6 explicit \texttt{application\_name}
overrides for cases where the category is empirically wrong for
scheduling purposes: Steam, EpicGames, Blizzard, and Electronic Arts are
tagged category ``Game'' by nDPI but are dominated by large client/patch
downloads (\(\rightarrow\) Bulk-Download), while TikTok and Facebook/
Instagram Reels are tagged ``SocialNetwork'' but carry continuous
short-form video (\(\rightarrow\) Video-Streaming). This taxonomy is
fully enumerated in \texttt{common/qos\_mapping.py} in the accompanying
code. The resulting class distribution is heavily imbalanced --
Web-Browsing 57.9\%, Delay-Sensitive 20.7\%, Video-Streaming 12.6\%,
Bulk-Download 8.8\% -- reflecting the natural composition of general
network traffic, which we address with class-balanced training weights
rather than resampling.

\subsection{Feature Engineering}

For each flow we extract exactly the first 10 packets' sizes, directions,
and inter-arrival times from the SPLT sequence NFStream already records
(which is itself prefix-truncatable: capturing with
\texttt{splt\_analysis=25} and using only the first 10 values is
equivalent to having captured with \texttt{splt\_analysis=10}). We fold
direction into signed packet size (positive for the flow-initiating
direction, negative for the reverse) rather than using a separate
direction feature, giving a 20-dimensional vector
\((\text{ps}_1,\ldots,\text{ps}_{10}, \text{iat}_1,\ldots,\text{iat}_{10})\)
per flow. This feature-extraction logic
(\texttt{common/splt\_features.py}) is imported by both the offline
training script and the real-time sniffer, so the two cannot silently
diverge.

\subsection{Model Training and Deployment Sizing}

We evaluate both a Random Forest and a LightGBM classifier
(\texttt{class\_weight=`balanced'}), split 80/20 using
\texttt{StratifiedGroupKFold} grouped by client IP address rather than a
plain random split -- since multiple flows from the same client can share
strong incidental correlations, a random split risks leaking client
identity into the test set and overstating accuracy. This yielded
348{,}173 training and 81{,}424 test flows spanning 2{,}411 and 615
distinct client IPs respectively.

An unconstrained Random Forest (100 estimators, no depth limit) reaches
0.886 macro F1 but pickles to 1.3\,GB, which we consider impractical for
a router or embedded-CPE deployment target. Bounding tree depth to 22 and
requiring at least 3 samples per leaf reduces this to 236\,MB at a cost
of 0.015 macro F1 (Table~\ref{tab:sizing}); we deploy this bounded model.
LightGBM underperformed the Random Forest in our setting (0.819 macro
F1) and was not selected.

\begin{table}[t]
\centering
\caption{Model size vs.\ accuracy trade-off (Random Forest).}
\label{tab:sizing}
\begin{tabular}{lccc}
\toprule
Configuration & Macro F1 & Pickle size \\
\midrule
Unbounded depth, 200 trees & 0.886 & 1.3\,GB \\
\(\text{max\_depth}=22\), \(\text{min\_leaf}=3\), 100 trees & 0.871 & 236\,MB \\
\bottomrule
\end{tabular}
\end{table}

\subsection{Real-Time Router Architecture}

The deployed sniffer tracks each observed flow by a direction-agnostic 5-tuple key. The first
packet seen for a key fixes which endpoint is the flow initiator,
matching NFStream's own directionality convention. On the 10th packet of
a previously unseen flow, the sniffer builds the 20-dimensional feature
vector and invokes the trained classifier; the resulting label is used to
install a persistent \texttt{iptables} mangle-table rule that sets a
firewall mark for both directions of that 5-tuple. A one-time \texttt{tc}
setup step creates 4 HTB leaf classes under one root qdisc on the
router's egress interface, each with an SFQ (Stochastic Fairness Queuing)
leaf so that multiple concurrent flows within one QoS class are scheduled
fairly rather than one flow starving its siblings, and 4 \texttt{tc}
filters route packets by firewall mark into the matching class. Because
enforcement is ``classify once per flow, then let the kernel forward
according to a static mark,'' steady-state per-packet cost after the
10th packet is a single netfilter table walk, not a repeated ML
inference.

Class rates and priorities are configured as fractions of the operator-
supplied link capacity rather than fixed absolute values, so the same
hierarchy is meaningful at any link speed: Delay-Sensitive receives the
highest scheduling priority (served first, bounding its queuing delay)
with a 30\% guaranteed-rate share; Video-Streaming receives the largest
guaranteed share (40\%) at the next priority level; Web-Browsing is the
default best-effort class (20\%); and Bulk-Download receives the smallest
guaranteed share (10\%) at the lowest priority, though it may still burst
up to a higher ceiling when the link is otherwise idle.

\subsection{Testbed}

We designed a Containernet topology (4 Docker client containers, one per
QoS class, connected via an OVS switch to a router container, which
connects via a second, bandwidth-constrained link to a server container)
so that all traffic shares one bottleneck link the router's \texttt{tc}
hierarchy controls -- exactly mirroring a home router's constrained WAN
uplink. This topology, plus a D-ITG traffic-profile generator and an
orchestration script that runs a ``baseline'' (classify only, flat
rate-limit, no per-class priority) phase followed by a ``QoS-enabled''
phase with the full mechanism active, is included in the accompanying
code (\texttt{testbed/topology.py}, \texttt{testbed/run\_experiment.py}).
Running it requires Containernet (a Mininet fork with Docker-backed
hosts) installed with root privileges for network-namespace and virtual-
interface management.

The results reported in Section~\ref{sec:results} were collected on a
functionally equivalent but topologically simplified pilot: two isolated
Docker bridge networks (LAN and WAN) joined by a dual-homed router
container performing real IP forwarding, with the router's
WAN-facing interface as the bottleneck the exact same \texttt{tc}/
\texttt{iptables} code shapes. We built this pilot because host root was
not available in our execution environment to run Containernet directly;
we verified real IP forwarding through the router container (ICMP TTL
decrementing by exactly 1 in both directions) and confirmed by direct
inspection (\texttt{tc class show}, \texttt{iptables -t mangle -L}) that
the installed HTB hierarchy and firewall-mark rules exactly match what
the Containernet topology would produce. The pilot and the full topology
share the router and traffic-generator container images
(\texttt{testbed/docker/*.Dockerfile}) and the identical
\texttt{router/sniffer.py} and \texttt{router/tc\_manager.py} code.

\subsection{Traffic Generation}

For each QoS class we define a D-ITG profile chosen to resemble that
class's traffic shape: Delay-Sensitive is small-packet, fixed-rate UDP
(160\,B at 50\,pkt/s, a VoIP-like codec rate); Video-Streaming is
large-packet, high-rate UDP (1400\,B at 400\,pkt/s); Bulk-Download is
large-packet, high-rate TCP (1400\,B at 900\,pkt/s, throttled in practice
by TCP congestion control and the bottleneck link); and Web-Browsing is
medium-packet TCP (800\,B at 60\,pkt/s). All four profiles run
concurrently against a shared 5\,Mbit/s bottleneck -- deliberately far
below the combined offered load -- so that the QoS mechanism has genuine
contention to resolve.

\subsection{Real-Traffic Replay and Independent Bandwidth Validation}
\label{sec:tcpreplay-method}

To directly test the domain-gap hypothesis in
Section~\ref{sec:results}\,D rather than only state it, we additionally
replay genuine packet captures through the live classifier using
\texttt{tcpreplay}. We take two pcaps from the tutorial's own data
collection module (1{,}459 and 7{,}794 packets respectively, containing
real DNS queries and real TLS/HTTPS connections to a mix of destinations),
rewrite their addresses with \texttt{tcprewrite} (\texttt{--srcipmap}/
\texttt{--dstipmap} pseudo-NAT plus \texttt{--enet-dmac} to target the
router's LAN-facing MAC and \texttt{--fixcsum} to repair checksums after
rewriting) so they route through the same pilot topology as
Section~\ref{sec:qos-pilot}, and replay them at topspeed from a client
container while the live sniffer classifies in real time -- this is
qualitatively different from the D-ITG case in that the packets are not
synthetically generated at all; they are unmodified real application
traffic, only relabeled at the IP layer to fit the testbed's addressing.

We additionally use \texttt{iperf3} as a bandwidth-measurement instrument
independent of both D-ITG and our own \texttt{tc} code, to (a) confirm the
pilot topology's baseline capacity is effectively unconstrained absent
deliberate shaping, and (b) cross-check that our HTB configuration
actually delivers the rate it claims to.

\section{Results}
\label{sec:results}

\subsection{Classification Performance}

Table~\ref{tab:classification} reports per-class precision, recall, and
F1 on the held-out, client-IP-disjoint test set (81{,}424 flows) for the
deployed (depth-bounded) Random Forest. Delay-Sensitive and Web-Browsing,
the two most common classes, are also the most accurately classified
(F1 0.941 and 0.930); Video-Streaming and Bulk-Download, the two least
common classes, are harder (F1 0.800 and 0.811), consistent with class
imbalance despite balanced training weights. Table~\ref{tab:confusion}
shows the confusion matrix: the dominant confusions are Bulk-Download and
Video-Streaming flows misclassified as Web-Browsing, plausible given that
many bulk/video transfers begin with a browser-mediated TLS handshake
before their bulk phase, which the first 10 packets may still be within.

\begin{table}[t]
\centering
\caption{Per-class classification performance (Random Forest, held-out test set, \(n=81{,}424\)).}
\label{tab:classification}
\begin{tabular}{lcccc}
\toprule
Class & Precision & Recall & F1 & Support \\
\midrule
Delay-Sensitive & 0.948 & 0.935 & 0.941 & 15{,}957 \\
Web-Browsing & 0.943 & 0.917 & 0.930 & 50{,}857 \\
Bulk-Download & 0.796 & 0.828 & 0.811 & 5{,}102 \\
Video-Streaming & 0.749 & 0.859 & 0.800 & 9{,}508 \\
\midrule
Macro avg & 0.859 & 0.885 & 0.871 & 81{,}424 \\
Weighted avg & 0.912 & 0.908 & 0.909 & 81{,}424 \\
\bottomrule
\end{tabular}
\end{table}

\begin{table}[t]
\centering
\caption{Confusion matrix (rows: true class, columns: predicted class).}
\label{tab:confusion}
\small
\begin{tabular}{lcccc}
\toprule
 & Bulk & Delay & Video & Web \\
\midrule
Bulk-Download   & 4{,}222 & 51 & 89 & 740 \\
Delay-Sensitive & 123 & 14{,}916 & 63 & 855 \\
Video-Streaming & 58 & 36 & 8{,}172 & 1{,}242 \\
Web-Browsing    & 902 & 739 & 2{,}591 & 46{,}625 \\
\bottomrule
\end{tabular}
\end{table}

The five most important features by Gini importance are the signed
packet size at positions 4, 8, 1, 5, and 10 of the 10-packet window,
together accounting for over a third of total importance -- i.e., no
single packet dominates; the classifier relies on the size trajectory
across the handshake rather than any one packet.

\subsection{Router Processing Latency}

Table~\ref{tab:latency} reports wall-clock timing from the live sniffer
across 8 real classification events captured during the pilot run
(Section~\ref{sec:qos-pilot}), broken into feature extraction, model
inference, and \texttt{tc}/\texttt{iptables} rule installation. Total
per-flow processing latency averaged 28.9\,ms (median 27.8\,ms, maximum
35.3\,ms). We caution that this was measured on a shared, virtualized
16-core development workstation running numerous unrelated processes
concurrently, not a dedicated router CPU; the measurement demonstrates
feasibility (processing completes in tens of milliseconds, negligible
relative to a flow's total duration) rather than a hardware-optimized
benchmark. Model inference dominates the total (24.1\,ms of 28.9\,ms);
feature extraction itself is negligible (0.06\,ms).

\begin{table}[t]
\centering
\caption{Router processing latency, live sniffer (\(n=8\)).}
\label{tab:latency}
\begin{tabular}{lccc}
\toprule
Stage & Mean (ms) & & \\
\midrule
Feature extraction & 0.059 & & \\
Model inference & 24.109 & & \\
\texttt{tc}/\texttt{iptables} enforcement & 4.763 & & \\
\midrule
Total (mean / median / max) & 28.931 & 27.821 & 35.343 \\
\bottomrule
\end{tabular}
\end{table}

\subsection{QoS Enforcement Effectiveness}
\label{sec:qos-pilot}

To isolate the enforcement mechanism's effect from classifier accuracy on
the (as discussed next) out-of-domain synthetic traffic, we drive
\texttt{tc}/\texttt{iptables} enforcement in this experiment from the
D-ITG flows' known ground-truth port rather than the live classifier's
prediction -- an oracle-labeling control that isolates ``does the
scheduler deliver differentiated treatment when given a correct label''
from ``does the classifier produce a correct label on this traffic.''
The classifier's own accuracy is independently established on real
traffic in Section~\ref{sec:results}\,A. We compare a flat, single-queue
5\,Mbit/s rate limit (baseline: no per-class differentiation, modeling a
QoS-unaware router) against the full 4-class HTB hierarchy at the same
5\,Mbit/s total capacity (QoS-enabled), running the identical 4
concurrent traffic profiles for both.

\begin{table}[t]
\centering
\caption{Per-class delay before vs.\ after enabling QoS enforcement (5\,Mbit/s bottleneck, single run).}
\label{tab:qos}
\small
\begin{tabular}{lrrr}
\toprule
Class & Before (ms) & After (ms) & Change \\
\midrule
Delay-Sensitive & 1{,}532.807 & 0.095 & $-$100.0\% \\
Web-Browsing & 2{,}590.788 & 0.144 & $-$100.0\% \\
Video-Streaming & 1{,}566.076 & 386.488 & $-$75.3\% \\
Bulk-Download & 6{,}113.880 & 4{,}062.180 & $-$33.6\% \\
\bottomrule
\end{tabular}
\end{table}

Table~\ref{tab:qos} shows the resulting average one-way delay per class.
Delay-Sensitive and Web-Browsing -- the two highest-priority classes in
our scheme -- see their average delay collapse from over a second to
sub-millisecond, a direct consequence of HTB's strict priority ordering
giving them first claim on the link whenever they have data queued.
Video-Streaming and Bulk-Download, the two lower-priority classes, still
improve relative to the chaotic flat-queue baseline (their delay drops
too, since the previously undifferentiated single queue is now at least
partitioned into fair per-class SFQ queues) but remain the most congested
classes, as intended: the whole point of the priority scheme is to
protect latency-critical traffic \emph{at the expense of} classes that
tolerate delay.

This trade-off is visible in throughput and loss, not just delay.
Bulk-Download's average bitrate fell from 1{,}558.6 to 966.8\,kbit/s under
enforcement -- it is deliberately the lowest-priority, smallest-
guaranteed-share class, so ceding bandwidth under contention is the
\emph{correct} behavior, not a defect. Video-Streaming's packet loss rose
from 6.77\% to 19.13\% under enforcement even as its delay improved: its
guaranteed 40\% share of a 5\,Mbit/s link (2\,Mbit/s) is still below its
offered load (\textasciitilde4.3\,Mbit/s unconstrained), so once its
queue is isolated from the other three classes it experiences more
concentrated loss than it did while occasionally opportunistically
sharing the undifferentiated baseline queue. We report this honestly
rather than only the favorable numbers: a real QoS deployment would size
each class's guaranteed rate against expected demand, and 5\,Mbit/s total
against four concurrent high-bitrate synthetic flows is a deliberately
severe stress case chosen to make contention visible within a short test
run, not a capacity-planning recommendation.

\subsection{Finding: Synthetic Traffic Does Not Resemble the Classifier's Training Domain}

Running the \emph{live} classifier (rather than the oracle labels above)
against the same D-ITG traffic produced a striking result: the four
D-ITG data flows were classified as Web-Browsing in 6 of 8 observed
instances and Video-Streaming in 1 of 8, essentially never matching their
intended class, while D-ITG's own signaling control channel (a short,
low-volume TCP connection ITGSend opens to ITGRecv before each transfer)
was classified Delay-Sensitive in all 8 instances. We attribute this to a
genuine domain gap: D-ITG generates statistically simple,
constant-inter-departure-time, constant-packet-size synthetic traffic,
whereas our classifier was trained on the first-10-packet SPLT signature
of real encrypted applications, whose early packets are dominated by a
TLS/QUIC handshake (variable-size ClientHello/ServerHello/certificate
exchange messages) rather than steady-state payload traffic. A
constant-rate synthetic stream simply does not contain that signature,
so the classifier -- reasonably, from its perspective -- falls back
toward the majority training class. D-ITG's control-channel connections,
being short and bursty, coincidentally resemble the interactive/
control-plane traffic (DNS, SNMP, remote access) that dominates our
Delay-Sensitive category.

This is a limitation of using D-ITG (or any purely synthetic generator)
to evaluate an SPLT-based early classifier end-to-end, not a defect in
the classifier itself, whose accuracy is independently and directly
measured on real traffic in Section~\ref{sec:results}\,A. We view
this as a useful negative result: end-to-end evaluation of early
classification systems should use real traffic captures (or realistic
replay, e.g.\ \texttt{tcpreplay} of real pcaps) for the classification
leg specifically, reserving synthetic generators like D-ITG for what they
are well suited to -- generating controlled, reproducible \emph{load} to
stress-test the enforcement mechanism once a flow's class is known. We
test this directly next, rather than leaving it as an untested claim.

\subsection{Confirmation: Real-Traffic Replay Produces Non-Degenerate Predictions}
\label{sec:tcpreplay-results}

Using the \texttt{tcpreplay} setup of Section~\ref{sec:tcpreplay-method},
we replayed two real pcaps (1{,}459 and 7{,}794 packets, 97 and 358
distinct flows respectively) through the live classifier. Of the 455
total real flows, only 15 reached the 10-packet threshold required for
classification at all -- itself a finding: real traffic contains many
short-lived flows (single DNS queries, brief HTTPS requests) that never
accumulate enough packets to classify, unlike D-ITG's flows, which are
long by construction. Of the 15 classified real flows, the live model
predicted Web-Browsing for 10 (real HTTP/HTTPS connections on ports 80
and 443) and Delay-Sensitive for 5 (including a UDP/443 flow, plausibly
QUIC, and several short TCP/443 connections) -- a plausible, varied split
given this capture is generic web/cloud traffic, and qualitatively
different from the D-ITG result: no single class dominates to the same
degree, and nothing here resembles the earlier near-total collapse onto
one class. We do not have independent ground-truth QoS labels for these
specific captured flows, so we report this as supporting evidence for the
domain-gap explanation rather than a second accuracy measurement, but the
contrast is exactly what that explanation predicts: real traffic, even
without known labels, produces classifier behavior that looks like
genuine discrimination rather than a default fallback.

Router processing latency on this real-traffic run (35--48\,ms per
flow) was consistent with the D-ITG-driven measurement in
Table~\ref{tab:latency}, confirming the latency figures are a property of
the classification pipeline itself rather than an artifact of synthetic
traffic's regular shape.

Separately, we used \texttt{iperf3} -- a bandwidth-measurement tool
independent of both D-ITG and our own \texttt{tc} code -- to sanity-check
the testbed. With no \texttt{tc} shaping active, an 8-second TCP transfer
between the same two containers reached 3.6\,Gbit/s (client-reported send
rate 10.9\,Gbit/s; the gap is normal send-buffer-vs-delivery behavior
under an effectively unconstrained virtual link), confirming that the
severe contention observed in Section~\ref{sec:qos-pilot} comes entirely
from our deliberately imposed 5\,Mbit/s cap and not from any incidental
limit of the Docker network itself. With the flat 5\,Mbit/s baseline
\texttt{tc} configuration active, the same test measured 4.77\,Mbit/s
received -- within 5\% of the configured cap -- independently confirming
that our HTB configuration enforces the rate it claims to, using a
measurement tool that shares no code with either the enforcement
mechanism or its evaluation.

\section{Limitations and Future Work}

\textbf{Root privilege requirement.} The full Containernet topology
requires host root for network-namespace and OVS management; our
reported measurements come from a topologically simplified but
mechanism-identical Docker pilot for this reason. Reproducing our results
on the full Containernet topology, or on physical hardware, is left as
direct future work using the code provided.

\textbf{Synthetic traffic domain gap.} As discussed above, D-ITG traffic
does not exercise the classifier realistically. We partially addressed
this with a \texttt{tcpreplay}-based real-traffic-replay experiment
(Section~\ref{sec:tcpreplay-results}), which showed qualitatively
healthier, non-degenerate classifier behavior on genuine traffic; a
remaining gap is that we lack independent ground-truth QoS labels for the
specific replayed captures, so this evidence is corroborating rather than
a second accuracy number. Future work should replay real traffic with
known labels (e.g.\ matched against the original nDPI-labeled dataset)
through the full enforcement pipeline to obtain that measurement.

\textbf{Single-run measurements.} The QoS before/after numbers in
Table~\ref{tab:qos} come from one 15-second run per condition rather than
repeated trials with reported variance, appropriate for demonstrating
that the mechanism functions correctly but not for precise effect-size
estimation. The router latency sample (\(n=8\)) is similarly small.

\textbf{Model size vs.\ depth trade-off.} We selected \(\text{max\_depth}=22\)
by inspection of a small sweep (Table~\ref{tab:sizing}); a systematic
Pareto study across depth, leaf size, and estimator count would likely
find a better-sized operating point, particularly for constrained
embedded-CPE targets rather than the general-purpose router we assume.

\textbf{Encrypted control-plane confusions.} The dominant confusion
classes (bulk/video traffic misclassified as web browsing) suggest the
first 10 packets sometimes still lie within a shared handshake phase
common to many application types; extending the feature window
adaptively, or combining SPLT with TLS/QUIC handshake metadata (SNI,
cipher suite, ALPN) where legally and contractually permissible, is a
natural extension.

\section{Conclusion}

We built and empirically validated an end-to-end pipeline from real
encrypted-traffic flow records to a deployable, size-constrained
classifier, a real-time router sniffer sharing its feature-extraction
code with the training pipeline, and a Linux \texttt{tc}/\texttt{iptables}
enforcement mechanism verified against a router-in-the-middle testbed.
The classifier reaches 0.871 macro F1 on 429{,}597 real flows at a
deployment-practical 236\,MB; the enforcement mechanism, evaluated
against ground-truth labels to isolate it from classifier accuracy,
reduces average delay for our two highest-priority classes from over a
second to sub-millisecond under a deliberately severe bottleneck. We also
report, rather than obscure, a real limitation: synthetic D-ITG traffic
does not resemble the real applications the classifier was trained on,
which is itself an actionable finding for how early-classification
systems should be evaluated end-to-end. We tested this directly with
\texttt{tcpreplay}-based replay of genuine packet captures, which produced
qualitatively healthier classifier behavior than synthetic traffic, and
cross-validated our \texttt{tc} enforcement mechanism against
\texttt{iperf3}, an independent tool sharing no code with our own
pipeline. All code, the trained model, and the raw experimental logs
underlying every number in this paper are available in the accompanying
repository.

\bibliographystyle{plain}
\bibliography{references}

\end{document}
